% !TeX program = xelatex
\documentclass{techpaper2026}

% 可选参数是页眉短标题；题目较长时建议填写。
\papertitle[二甲醚清洁燃料均质压燃燃烧数值模拟研究]
  {二甲醚清洁燃料均质压燃燃烧数值模拟研究}
\paperauthor{XXX}
\paperstudentid{XXX}
\paperclass{XXX}

\begin{document}

\makepapertitle

\begin{paperabstract}
均质充量压缩着火（HCCI）燃烧是一种可实现高效低污染的燃烧方式。
本文以二甲醚（DME）为例，采用化学动力学模型研究压缩比、燃空当量比、
进气温度、发动机转速及废气再循环等因素对着火和燃烧过程的影响。
结果表明，DME 的 HCCI 燃烧具有明显的低温和高温两阶段放热特征。
此处替换为论文摘要，摘要正文使用中文五号宋体、英文五号 Times New Roman，
首行缩进两个汉字并采用单倍行距。
\end{paperabstract}

\paperkeywords{均质充量压缩着火，化学动力学，数值模拟，二甲醚，EGR}

\section{引言}

HCCI 发动机的着火与燃烧过程受到燃料化学反应动力学的显著影响。
中文正文为五号宋体，English text is set in 10.5,pt Times New Roman，
每段首行缩进两个汉字，单倍行距。中文文献可写作\citep{li1959}，
两位作者的西文文献可写作\citep{hanahan2000}，多位作者可写作
\citep{you1980}；多篇同时引用可写作\citep{hanahan2000,you1980}。

\subsection{HCCI 的数值模拟研究现状}

本节用于展示二级标题格式。二级标题顶格书写，序号后空一个汉字，
采用小四号黑体和单倍行距。

\subsubsection{HCCI 数值模拟模型}

三级标题和四级标题均缩进两个汉字，采用小四号宋体。

\paragraph{单区模型}

此处填写四级标题下的正文内容。对总项包含的分项可使用以下编号：

\begin{enumerate}
  \item 单区模型；
  \item 双区和多区模型；
  \item 多维模型。
\end{enumerate}

\section{材料与方法}

\subsection{计算方法}

正文中的公式应另起一行并按全文连续编号。较长公式优先在等号处换行，
否则在加、减、乘、除等运算符处换行，并将运算符写在新行行首：

\begin{equation}
  m=\sum_{k=1}^{K}m_k .
\end{equation}

\begin{align}
  L(Z,S,U,E,F)
    ={}& \sum_{i,j}^{n}\lvert x_i-x_j\rvert_2^2s_{ij}
       +\lambda_1\lVert U\rVert_*
       +\lambda_2\lVert E\rVert_1 \notag\\
     &+2\lambda_3\operatorname{Tr}(F^{\mathsf T}L_SF)
       +\langle C_1,X-XZ-E\rangle \notag\\
     &+\frac{\mu}{2}\left(
       \lVert X-XZ-E\rVert_F^2+
       \lVert Z-S\rVert_F^2+
       \lVert Z-U\rVert_F^2\right).
\end{align}

\section{结果}

\subsection{热力学性质}

表题置于表格上方，表序在左、表题在右，中间空一格，末尾不加标点。

\begin{table}[htbp]
  \centering
  \caption{选取组分的热力学性质}
  \label{tab:properties}
  \begin{tabular}{cccc}
    \toprule
    组分 & $H_f$ (kcal/mol) & $S_f$ (kcal/mol) & $C_p$ (kcal/mol) \\
    \midrule
    A1 & 100 & 100 & 100 \\
    A2 & 110 & 105 & 102 \\
    A3 & 120 & 108 & 104 \\
    \bottomrule
  \end{tabular}
\end{table}

% 跨页表格示例：续页重复表头，并在右上角显示“续表 n”。
\begin{longtable}{>{\centering\arraybackslash}p{0.18\textwidth}
                  >{\centering\arraybackslash}p{0.24\textwidth}
                  >{\centering\arraybackslash}p{0.24\textwidth}}
  \caption{跨页表格示例}\label{tab:long}\\
  \toprule
  方案 & 参数一 & 参数二 \\
  \midrule
  \endfirsthead
  \multicolumn{3}{r}{\zihao{5}\PaperSong 续表~\ref{tab:long}}\\
  \toprule
  方案 & 参数一 & 参数二 \\
  \midrule
  \endhead
  \bottomrule
  \endfoot
  A & 1.0 & 2.0 \\
  B & 1.1 & 2.1 \\
  C & 1.2 & 2.2 \\
\end{longtable}

\subsection{结果曲线}

每幅插图应有连续的图序和图题。将下方占位框替换为
\texttt{\textbackslash includegraphics} 命令即可插入图片。

\begin{figure}[htbp]
  \centering
  \fbox{\rule{0pt}{45mm}\rule{0.72\textwidth}{0pt}}
  \caption{气缸压力随曲轴转角变化的曲线}
  \label{fig:pressure}
\end{figure}

\section{讨论}

此处讨论研究结果、误差来源、适用范围及与既有研究的差异。
公式、表格和图片应分别按全文统一连续编号，不得重复或跳号。

% 参考文献须手工按“中文在前、外文在后”排序：
% 中文按第一作者姓名拼音排序，西文按第一作者姓氏字母排序。
\begin{thebibliography}{99}

\bibitem[李薰(1959)]{li1959}
\JournalRef{李薰}{十年来中国冶金科学技术的发展}{金属学报}{1959}{4}{3--10}

\bibitem[谭丙煜(1982)]{tan1982}
\BookRef{谭丙煜}{怎样撰写科学论文}{}{沈阳}{辽宁人民出版社}{1982}{59--76}

\bibitem[王磊(2014)]{wang2014}
\ThesisRef{王磊}{中国人非小细胞肺癌 PIK3CA 基因突变及 FGFR 融合基因研究}
  {博士学位论文}{上海}{复旦大学}{2014}{103--115}

\bibitem[Eissen(1974)]{eissen1974}
\BookRef{Eissen H. N.}{Immunology: an introduction to molecular and cellular principles
  of the immune responses}{5th ed}{New York}{Harper and Row}{1974}{40--49}

\bibitem[Hanahan \& Weinberg(2000)]{hanahan2000}
\JournalRef{Hanahan D., \& Weinberg R. A.}{The hallmarks of cancer}
  {Cell}{2000}{100}{57--70}

\bibitem[You et al.(1980)]{you1980}
\JournalRef{You C. H., Lee K. Y., Chey R. F. et al.}
  {Electrogastrographic study of patients with unexplained nausea, bloating and vomiting}
  {Gastroenterology}{1980}{79}{311--314}

\end{thebibliography}

\end{document}
